\section{A complete balanced three-stack realization}\label{sec:bb-balanced-three-stack}

In a positional code with base $B$, replacing a prefix of length $r$
by one of length $s$ gives an affine map with slope $B^{r-s}$.
Thus a three-stack instruction has determinant one when it preserves
the sum of the three prefix lengths.  A third stack supplies reversible
history under this constraint: we first move a fresh blank from beyond
a work-tape boundary onto the history stack, then replace that blank
by the record.

\begin{theorem}[Balanced three-stack realization]
\label{thm:bb-balanced-three-stack}
For every machine and finite input, there are effective coefficients
$f_0,f_1$ such that $f=f_0+\nu f_1$ on $\R^3$ has an explicit
zero-data solution $(U,0)$.  Force and velocity have one compact
spatial support and bounded mixed derivatives of every order, and
are one-periodic after $t=1$.  At the fixed label $x_*=(4,0,0)$,
\[
 \exists t\ge0:\quad X(t;x_*)\in(-1,2)^3
 \quad\Longleftrightarrow\quad\text{the machine halts}.
\]
For a fixed machine its repeated field is independent of the input.
The comparison class is $u\in C H^2\cap C^1L^2$, bounded $u,\nabla u$
on finite intervals, and pressure modulo constants in $C H^1$.
The formula is valid for every fixed real $\nu>0$, with effective
evaluation when $\nu$ is computable and relative evaluation otherwise.
\end{theorem}
\begin{proof}
Normalize the finite machine to one halting state $h$, directing missing
nonhalting instructions to it by unchanged-letter stay moves. This preserves
initial halting as well. Let $\Gamma$ be the work alphabet and
$b\in\Gamma$ its blank. Use three infinite top-first stacks $L,R,J$.  A rule replaces three
finite prefixes, leaving their tails untouched.  Introduce boundary
$\star$, marked work letters $\bar a$, and controls $A_q,S_q,T_q,D_q$;
only $A_h$ is terminal.  For each $q\ne h$, the preparatory rules are
\begin{align}
 (A_q;\varepsilon,a,\varepsilon)&\to(S_q;\bar a,\varepsilon,\varepsilon),
 \label{eq:bb-balanced-entry}\\
 (S_q;\varepsilon,c,\varepsilon)&\to(S_q;c,\varepsilon,\varepsilon),
 \label{eq:bb-balanced-out-sweep}\\
 (S_q;\varepsilon,\star b,\varepsilon)&\to(T_q;\varepsilon,\star,b),
 \label{eq:bb-balanced-blank}\\
 (T_q;c,\varepsilon,\varepsilon)&\to(T_q;\varepsilon,c,\varepsilon),
 \label{eq:bb-balanced-return-sweep}\\
 (T_q;\bar a,\varepsilon,\varepsilon)&\to(D_q;\varepsilon,a,\varepsilon).
 \label{eq:bb-balanced-restore}
\end{align}
Here $a,c\in\Gamma$.  For $\delta(q,a)=(q',d,\sigma)$ add
\begin{equation}\label{eq:bb-balanced-step-table}
\begin{array}{c|c|c}
\sigma&\text{source}&\text{target}\\\hline
0&(D_q;\varepsilon,a,b)&(A_{q'};\varepsilon,d,\kappa)\\
-1&(D_q;c,a,b)&(A_{q'};\varepsilon,cd,\kappa)\quad(c\in\Gamma)\\
1&(D_q;\varepsilon,ac,b)&(A_{q'};d,c,\kappa)\quad(c\in\Gamma)\\
1&(D_q;\varepsilon,a\star b,b)&(A_{q'};d,b\star,\kappa).
\end{array}
\end{equation}
Every individual row occurrence, including its choice of $c$, receives
a distinct record symbol $\kappa$.

\emph{Full-domain separation and initialized behavior.}
\label{lem:bb-balanced-processor}\label{proof:bb-balanced-processor}
Every rule preserves total prefix length: it is one in the transfer
rows, two in the blank-supply row, and respectively two, three, three,
four in the final rows.  Source branches at $A_q$ test the right top;
at $S_q$ the right top separates work letters from the boundary;
at $T_q$ the left top separates work and marked letters.  At $D_q$
the right top fixes the work instruction, with the next right symbol
or left top separating its remaining branches.  Target branches at
$S_q$ have distinct marked or ordinary left tops; at $T_q$ distinct
boundary or ordinary right tops; at $D_q$ distinct restored $a$;
and at an $A$ control distinct history tops $\kappa$.
This proves separation for arbitrary independent tails.

Initialize
\[
 (A_{q_0};b^\infty,v\star b^\infty,b^\infty),\qquad
 v=\begin{cases}\omega,&\omega\ne\varepsilon,\\b,&\omega=\varepsilon.
 \end{cases}
\]
At a checkpoint, $L$ lists cells strictly left of the head and
$R=a_1\cdots a_m\star b^\infty$, $m\ge1$, lists the current cell
and a finite right block; beyond it the tape is blank.  $J$ contains
completed records, newest first.  The first rule marks $a_1$ while
moving it left; $m-1$ forward transfers expose $\star$.  The blank
row removes one filler blank beyond that boundary and pushes it on
$J$, without changing the infinite blank filler.  The reverse transfers
and restoration of $a_1$ restore exactly the original two work stacks.
The final rule replaces the fresh history blank by its record and
performs the work instruction.  A left move pops the next cell from
$L$; a right move with $m\ge2$ pops the current cell from $R$;
when $m=1$ the last row exposes a blank and retains a nonempty right
block.  Thus the invariant holds and one work step costs
$1+(m-1)+1+(m-1)+1+1=2m+2$ rules.  An initially halted input needs
none.  The records recover the head displacements as well.

\emph{From balanced words to solid boxes.}
Give the enlarged alphabet $\Sigma$ the odd digits
$1,3,\ldots,2|\Sigma|-1$ in base $B=2|\Sigma|+1$.
For a finite word $v$ put $e(v)=\sum_{j=1}^{|v|}d(v_j)B^{-j}$;
for an infinite stack $\xi$ put
$e(\xi)=\sum_{j\ge1}d(\xi_j)B^{-j}$.
The prefix interval is $I(v)=e(v)+[0,B^{-|v|}]$, with
$I(\varepsilon)=[0,1]$.  Every infinite tail satisfies
\[
 \frac1{B-1}\le e(\xi)\le\frac{2|\Sigma|-1}{B-1}<1,
\]
so its code lies strictly inside each applicable prefix interval.
Distinct incompatible prefixes have positive gaps, since digits differ
by at least two while the remaining interval has width at most one unit
in the differing place. Place control $s$ at
$o_s=(4i_s,0,0)$, with $i_{A_h}=0$ and distinct positive integers
for the others.  The code of $(s;L,R,J)$ is
$o_s+(e(L),e(R),e(J))$.
A rule with input prefixes $\alpha_1,\alpha_2,\alpha_3$
and output prefixes $\alpha'_1,\alpha'_2,\alpha'_3$ gives boxes
\[
 Q_i=o_s+\prod_jI(\alpha_j),\qquad
 Q_i'=o_{s'}+\prod_jI(\alpha'_j).
\]
Both families are separately disjoint by the preceding full-prefix
separation.  In control-relative coordinates their exact maps are
\begin{equation}\label{eq:bb-balanced-affine-rule}
 y_j=e(\alpha'_j)+B^{|\alpha_j|-|\alpha'_j|}
                    (x_j-e(\alpha_j)).
\end{equation}
The identity $e(\alpha\xi)=e(\alpha)+B^{-|\alpha|}e(\xi)$
proves their action on every code.  Initialized stacks are eventually
blank, with rational codes
$e(vb^\infty)=e(v)+B^{-|v|}d(b)/(B-1)$.
The affine determinant equals
$B^{\sum|\alpha_j|-\sum|\alpha'_j|}=1$, exactly by balance.
Apply Theorem~\ref{lem:bb-balanced-box-routing} to these full boxes.

\emph{Every intermediate time.}
All half-widths are at most $1/2$, so that lemma uses $R=1$.
Sources have first coordinate at least four, since no rule leaves
$A_h$.  Targets with nonhalt control have the same property.
For such a branch the initial expansion and final contraction preserve
first coordinate at least four.  During deformation its center has
first coordinate at least $4\Lambda$.  Every scaled halt target
center has first coordinate at most $\Lambda$, whereas the endpoints
of each nonhalt parking path have first coordinate at least
$4\Lambda$.  Since $\Lambda+4<4\Lambda$, the original segment
cannot meet an obstacle cube around a halt target.  Any encountered
obstacle has first coordinate at least $4\Lambda$, so its boundary
detour remains at least $4\Lambda-4$.  Subtracting a half-width at
most one still leaves first coordinate greater than two.  Hence
\begin{equation}\label{eq:bb-balanced-no-premature-event}
 \Phi_\tau(Q_i)\subset\{x_1>2\}\quad(0\le\tau\le1)
 \quad\text{for every branch with nonhalt target}.
\end{equation}
This includes delivery after halt targets have already been filled.

Finally load $x_*=(4,0,0)$ along the straight segment to the rational
initial code using the translation part of Lemma~\ref{lem:bb-curl}.
It avoids $(-1,2)^3$ for a nonhalt initial control.  Repeat the box
field from time one.  Induction at integer times applies the exact
full-box maps; the finite compiler cycles reach a code in $(0,1)^3$
precisely upon halting.  Equation~\eqref{eq:bb-balanced-no-premature-event}
excludes every intermediate false visit.  The residual proposition
gives the force and uniqueness; compact support gives all finite-slab
Sobolev conditions and uniformly bounded kinetic energy.  Every choice
was finite and effective, including the obstacle detours.  The input
enters only the first loading path, completing the proof.
\end{proof}
